\label{sec:heine-borel}

\begin{topicbox}{The Capstone of the Real Line}
\begin{exposition}
Heine--Borel closes the chapter by proving the converse direction: on $\mathbb{R}$, closed and bounded is enough for compactness. The keystone is that every closed bounded interval $[a,b]$ is compact; the general statement follows because a closed subset of a compact set is compact.
\end{exposition}
\end{topicbox}
\begin{theorembox}{Theorem (Closed Bounded Intervals Are Compact)}
\begin{theorem}[Closed Bounded Intervals Are Compact]
\label{thm:closed-bounded-interval-compact}
Every closed bounded interval $[a,b]\subseteq\mathbb{R}$ is compact.
\hyperref[prf:closed-bounded-interval-compact]{Go to proof.}
\end{theorem}
\end{theorembox}
\begin{remark*}[Standard quantified statement]
\[
\forall a,b\in\mathbb{R}\;\bigl(a\le b \Longrightarrow [a,b]\text{ is compact}\bigr).
\]
\end{remark*}
\begin{remark*}[Interpretation]
This is the hard half. The standard proof uses completeness directly: let $S$ be the set of $x\in[a,b]$ such that $[a,x]$ has a finite subcover, and show $\sup S=b$ is attained. Bisection (the nested-interval method) gives the same result. Either route shows completeness is what makes the line compact in the large.
\end{remark*}
\begin{dependencies}
\begin{itemize}
  \item \hyperref[def:compact-set]{Compact Set in $\mathbb{R}$}
  \item \hyperref[def:closed-interval]{Closed Interval}
  \item \hyperref[thm:lub-property-r]{Least-Upper-Bound Property}
\end{itemize}
\end{dependencies}
\begin{theorembox}{Theorem (Heine--Borel Theorem for $\mathbb{R}$)}
\begin{theorem}[Heine--Borel Theorem for $\mathbb{R}$]
\label{thm:heine-borel}
A set $K\subseteq\mathbb{R}$ is compact if and only if $K$ is closed and bounded.
\hyperref[prf:heine-borel]{Go to proof.}
\end{theorem}
\end{theorembox}
\begin{remark*}[Standard quantified statement]
\[
K \text{ compact} \Longleftrightarrow (K \text{ closed} \land K \text{ bounded}).
\]
\end{remark*}
\begin{remark*}[Interpretation]
Heine--Borel is the bridge into Analysis~I: it certifies that the closed bounded sets --- the ones on which the Extreme Value and uniform-continuity theorems live --- are exactly the compact ones. It is stated here as the equivalence \emph{on $\mathbb{R}$}; the open-cover side is the definition that later carries to the torus, while ``closed and bounded'' does not.
\end{remark*}
\begin{remark*}[Exposition]
With Heine--Borel in hand the chain is complete: number construction $\to$ embedded number systems $\to$ the real line as an ordered metric space $\to$ Analysis~I. Completeness, introduced as the last rung of the embedding ladder, is exactly the property that makes the line compact in the large, and compactness is the hypothesis Analysis~I leans on.
\end{remark*}
\begin{dependencies}
\begin{itemize}
  \item \hyperref[thm:compact-implies-closed-bounded]{Compact Subsets of $\mathbb{R}$ Are Closed and Bounded}
  \item \hyperref[thm:closed-bounded-interval-compact]{Closed Bounded Intervals Are Compact}
\end{itemize}
\end{dependencies}
